\emph{Carried over verbatim from the full report (draft v2.4, 23 September 2026), which is released with this report. Section and appendix numbers in this appendix refer to that report; Appendix~E maps them to this one.}

\subsection*{Full report, Section 11: Corrections since draft v1.0}

Appendix D lists fifty-four corrections to statements of our earlier drafts and dated files, each made in place in this draft. Items 1 to 4 were listed in draft v2.0 (14 September), which draft v2.1 overwrote and which is not preserved as a separate file. Items 5 to 18 come from an adversarial verification of draft v2.1 against the dated files and the correspondence (21 September). Items 19 to 53 come from a second adversarial verification, of draft v2.3 (23 September), which checked every factual sentence again against the dated files, the saved outputs and the public records of Codabench, GitHub, Hugging Face and the competition website, recomputed every figure it could, and also extended items 1, 4, 6, 8, 10 to 14, 16 and 17 with corrections of its own. Item 54 comes from a third adversarial verification, of this draft before its release (23 September), which also revised items 1, 4 to 6, 11 to 22, 29 to 32, 34 to 39, 41 to 43, 46 and 49 to 53, adding to several of them corrections of statements of earlier drafts (for example items 18, 31, 42, 43, 49, 51 and 52); its other findings, on sentences of this draft and on the dated notes, were corrected before release. A fourth, narrower check of the changes made after the third found further problems of attribution and wording, which revised items 1, 18, 21, 22, 34, 42, 51, 53 and 54 and the dated notes; they too were corrected before release. The scripts the verifications required were saved on 21 and 23 September (Section 13).

The corrections that change a conclusion include these: the classifier behind H10 was not the official baseline (item 24); the H13 ceiling came from thresholds chosen in-sample, not from a constant (item 25); the E11 analysis applied no multiplicity correction, and H30 would have failed Holm over its two tests (item 30); 26 of the 164 DeepSeek labels used as ground truth were derived by us (item 33); the organizers' DeepSeek labels come from a sampling regime of about ten answers per variant, not one (item 34); the organizers' probe artifact orders our labels, under its own inference aggregation, mostly through problem length (0.57 once residualized on log character count), as its own control probes trained on shuffled labels do (item 46); the plug-in simulation overstates reliability on our own labels by about 0.4 (item 41); on Olmo the seed accounts for about half of the instability under either rule, so which variants are drawn does not dominate there (items 39 and 43); the re-execution drift makes the official score channel's ordering irreproducible between identical runs but is not shown to change whether the channel carries signal (item 45); the cause of the numerical states is not identified (item 44); the budget guard is not shown to prevent timeouts (item 23); and the gap to the official DeepSeek labels is not only one of severity (item 43).

Our claim now is narrower and measured: apart from rule T, which exploits a subpopulation absent from \texttt{val-sample}, none of the channels available to us that we scored in accuracy, and not the organizers' own probe artifact, clears the majority constant by more than about one standard error on public ground truth. We make no claim about methods we have not seen. Apart from the public leaderboard figures cited in Sections 3.3 and 6, in the introduction of Appendix C and in Appendix D (items 1, 17 and 52), and predictions P1 and P2 in Appendix C, all of which concern scores and submissions, not methods, nothing in this report should be read as a statement about any participant.

\subsection*{Full report, Appendix D: Corrections to our earlier statements}

Each item names the section corrected and the drafts or dated files in which we found the statement; items 5 to 18 refer to draft v2.1 unless they name another. The correction is made in place in the section cited, and the dated files carry correction notes of 21 and 23 September.

\noindent\textbf{1.} Drafts v1.0 to v2.3 (Section 3.3) explained a flat leaderboard by the shipped baseline being a constant on unloadable models (v1.0: ``explains why the mid-warmup leaderboard clustered at exactly 0.53 and 0.47''; v2.1 to v2.3: ``which is why the mid-warmup leaderboard clustered at the two constant values''); our message of 30 August had described that constant behaviour on the warmup set only, without the leaderboard inference. The baseline falls back for lack of a shipped artifact, not of a cached model, and it scored 9/17 itself, so it could not account for the rows at 0.47; an explanation of other participants' scores by a presumed method is withdrawn (3.3). In the main phase the leaderboard moved: by 13 September individual submissions had reached 0.95 on Main and 1.00 on Small (recorded in PREREG\_V8). On 16 September both evaluation tasks were replaced (21 and 14 items). Until the morning of 21 September, before our uploads, no submission on the new tasks had exceeded 14/21 on Main or 9/14 on Small, against 10/21 and 7/14 for the all-robust constant (6.3); at 10:02 UTC that day a submission reached 11/14 on Small, and on 23 September at 03:40 UTC one reached 18/21 on Main. Under the declared rule the public table lists every scored submission as its own row (614 rows on 14 September, fifty per page, as recorded in PREREG\_V8), so the printed rank is per upload, not per participant. The narrower claim that replaces the explanation is stated in Section 11, with the exception of rule T that drafts v2.1 to v2.3 omitted there.
\par

\noindent\textbf{2.} During 8.1, after three identical repetitions, we concluded in an internal message that the drift came not from the GPU or the pipeline but from the shared environment. Repetitions four and five falsified that within the hour; the conclusion never left the project. The corrected finding is in 8.1.
\par

\noindent\textbf{3.} An internal analysis had stated that the shipped decision threshold ``costs 43 points'' because it predicted 66.5\% robust on our A100-extracted features. On the competition hardware it predicts 25.6\% robust in two of three identical extractions and 33.5\% in the third, against a 24.4\% rate, and costs 12 to 16 points against the constant, not 43. The 66.5\% came from our own re-implementation of the extraction (item 5); the claim was never published and is withdrawn here so that the record is complete.
\par

\noindent\textbf{4.} In an internal note of 13 September (PREREG\_V8, withdrawn there on 14 September) and in two paragraphs of PREREG\_E13 we inferred from the 0.41 test-retest of our regeneration that no method could exceed roughly 0.7 accuracy, and that leaderboard scores above it were therefore anomalous. The inference used our process's reliability where the organizers' was required, and the arithmetic was also wrong: with a reliability of 0.41 the bound on agreement is its square root, 0.64, which allows 0.82 accuracy on a balanced set, not 0.70. The inference is withdrawn.
\par

\noindent\textbf{5.} Draft v2.1 (Section 2.3(b)) reported that \texttt{generation\_frac\_selected\_is\_top1} is identically 1.0 at inference, so that one fifth of the forest routes every production sample the same way (the feature carries 17.1\% of the splits). That was true only of our own A100 re-implementation of the extraction; in the official extraction on the organizers' hardware the feature is never 1.0 and falls mostly inside the split range. The defect is withdrawn, with the half of R5 that depended on it.
\par

\noindent\textbf{6.} H2, the only warmup hypothesis about a property of the problem that had survived, survived only because the same 52 measurements, published under two model identifiers, were counted as two models, which made 24 problems multi-model by duplication. With the identifiers merged it dies by its own criterion (z = 0.88; $-$1.3 points against the per-model majority; Section 4), and also by E1's original test (z = 1.81 against 2). Drafts v2.2 and v2.3 also called it ``the only warmup signal that had survived''; H10, H16 and H18 had survived as well.
\par

\noindent\textbf{7.} Section 8.1 stated that different numerical states share none of the forty generations, that the mechanism is not per-item noise, and that contention is the evaluation workers' condition. Different states share few generations but not none (30 of the 91 pairs share at least one, one pair shares 9 and two batch-1 runs share 38), and the workers' condition is not something we could observe. It also stated that the two remedies make generation bit-reproducible; they did so in three sequential runs on one card and were not tested under contention or across cards.
\par

\noindent\textbf{8.} Section 8.2 described one of the three extractions as run ``with the determinism flags''; only the cuBLAS workspace variable was set, not the remedy of 8.1. It stated the orientation of the official channel backwards: the regressor predicts accuracy decay, so in the orientation of its own rule the two states give 0.487 and 0.602, not 0.513 ``in the expected direction'' and 0.398 against it. It called the shipped threshold well calibrated on this hardware (25.6\% predicted robust against a 24.4\% rate), which holds in one numerical state only; in the other it predicts 33.5\%. It compared these runs with an A100 figure (0.649) that came from our own extraction, was attributed to Section 5, which does not contain it, and filled a feature our extractor did not record with the exponential of the mean log-probability; with the mean top-1 probability, which the official extraction returns under greedy decoding, the figure is 0.64 (0.36 in the rule's orientation). And it did not state that E16 was not pre-registered.
\par

\noindent\textbf{9.} Section 7.2 attributed 91\% of DeepSeek's residual instability to the sampling seed alone; DeepSeek's A$'$ condition, through a queueing error, also used a different prompt, so that share was not measured. The pre-registered reading grid of E13, which classifies Qwen3.5-4B as unstable by construction, had been overridden without being declared (now declared in 7.3), and gpt-oss-120b, pre-registered for labeling, was not labeled.
\par

\noindent\textbf{10.} Section 7.3 described the organizers' variants as generated from symbolic reasoning chains and verified by mathematicians. The organizers' proposal (arXiv 2607.13899) does describe symbolic reasoning chains verified by mathematicians, but the published variants of the five families that 7.2 regenerates are attributed to language models in their \texttt{permutation\_source} field; the sentence conflated the two and is withdrawn, and 7.3 now states that the simulated value (0.94 on average, median 0.93) is optimistic, not evidence that the label is stable.
\par

\noindent\textbf{11.} Section 8.3 described the probe's permutation null as repeating the whole nested procedure; it used 60 permutations instead of the pre-registered 500 and a reduced regularization grid, and these deviations were not declared (re-run on 23 September, the reduced null reproduces its p95 of 0.574 and the full grid gives 0.570). It reported the organizers' probe artifact through an aggregation that differs from the artifact's own inference code (0.739 against 0.709), and said that the probe reads length ``from the lexical layers'', which no analysis shows. It also gave the probe's gain over character count from a single fold split (+0.008 in draft v2.1, about one point in v2.2); over 20 splits the gain is about two points (0.6 to 3.2).
\par

\noindent\textbf{12.} Draft v1.0 (Section 4.1) and draft v2.1 (Section 4) stated that grouped validation is what killed H1. Under an ungrouped 5-fold split, with the same code and data, H1 also dies ($-$0.9, $-$0.3 and $-$2.1 points over three seeds, against a bar of 2.7): grouping deepened the loss but did not cause it, and the note is withdrawn. Draft v2.1's count of dead hypotheses (nineteen) did not match its own table and is restated.
\par

\noindent\textbf{13.} Section 13 stated that no evaluation data was read, inferred or probed. Aggregate facts about the evaluation sets were logged and used; they are now listed there. Drafts v1.0 and v2.1 also stated in Section 3 that we logged only aggregate counts, timings and model\_id; from 31 August (v5, submission 909853) the logs also carried, per model batch, the minimum and maximum regressor score and the threshold applied, which on the two-item warmup batches are the two item scores. Drafts v2.1 to v2.3 omitted from the list in Section 13 the per-item token counts of 907967 and the fact that V8's length summaries give every item's length on batches of three items or fewer.
\par

\noindent\textbf{14.} On 1 September we wrote to the organizers, as Section 7 of draft v1.0 also stated, that two of the four warmup models produced identical per-item regressor scores across four runs and that the other two visited two and three distinct states. Our logs recorded only each model's minimum and maximum score: for gemma-3-27b-it and DeepSeek, with two items each, that is every item, but for gpt-oss-120b it is two of nine and for Olmo-3-7B-Instruct two of four. The per-item statement is withdrawn for gpt-oss-120b, and the two states reported for Olmo-3-7B-Instruct are a lower bound.
\par

\noindent\textbf{15.} On 30 August we wrote to the organizers, and Sections 6 and 8 of draft v1.0 repeated with a range of 30 to 40 items, that the official baseline would not finish beyond about 30 items, extrapolating from about 120 s per item on a batch of two. At the pipeline's batch size of 8 the measured cost is 26 s per item (8.1), and our five-model main-phase entry ran 21 items in 838 s.
\par

\noindent\textbf{16.} Section 6.3 stated that V8 showed that length carries nothing on the evaluation set, and the abstract that it scored exactly the majority constant there. The majority constant was never measured on that set, and one evaluation on 21 and 19 items cannot show that a feature carries nothing. Section 8.3 of drafts v2.1 to v2.3 kept a cross-reference to the withdrawn reading (``Section 6.3 measured what length is worth on the evaluation set''), and PREREG\_V8, PREREG\_V8b and PREREG\_E15 repeat it; it is withdrawn in their dated notes of 23 September.
\par

\noindent\textbf{17.} Section 6 of draft v1.0 and Section 9 of draft v2.1 attributed to the organizers an explanation of the cache gap (``an outdated configuration'') that they did not give, and Appendix C of draft v2.1 stated that the predictions were written before the main-phase leaderboard moved and left unchanged. The leaderboard had already moved: entries stood above the organizers' own entries on both tracks from 3 September, and on 6 September one participant's entries rose to five items above them on both tracks. P2 is anonymized.
\par

\noindent\textbf{18.} Appendix B of draft v2.1, which presented itself as unchanged from v1.0, dropped the review's verdict that v1.0 had stated; neither of these two drafts said that the entry's pre-registration forbade submitting it if the review refuted it, or that the author overrode that rule. Appendix B of drafts v2.2 and v2.3 said both, but opened by stating that v1.0 had not reported the review's verdict. Appendix B now corrects both.
\par

\noindent\textbf{19.} Dates. The warmup ended on 31 August 2026 (UTC; the Codabench phase record gives 15:00 UTC), not on 1 September as all drafts said. The organizers cached the four warmup models on 30 August, the day of our report, not on 31 August (all drafts: v1.0 in Section 6, which also dated our message 31 August; v2.1 in Section 9; v2.2 and v2.3 in Sections 3.2 and 9 and in Appendix B's ``a container change made that same day''). The rule on training data was dropped from the website's rules on 25 August and only confirmed in Discord on 5 September (v2.1 to v2.3, Sections 7 and 9; VERIFICA\_classifica\_2026-09-14). The rule of the main-phase labels was published with the train set on 16 September, not 17 (v2.1 to v2.3 and PREREG\_E13b). The simulation of 7.3 was first run on 13 September, not 14 (its note is misdated). And \texttt{permutation\_type} was dropped from the import script on 1 September by the organizers' commit that added the Small track, not on 2 September by PR \#5 (v2.1 to v2.3).
\par

\noindent\textbf{20.} Related work and scope (v2.2, v2.3). The first citation named no author and presented as already documented the variation at a fixed configuration that 8.1 measures, while the cited post traces nondeterminism to kernels that are not invariant to batch size; the second used the title of the paper's first arXiv version. Section 1 also described 8.1 and 8.2 as measured ``inside this benchmark's own evaluation environment'', where only H18 ran inside the evaluation container, 8.1 and 8.2 ran on a VM with the workers' GPU model and image, and no measurement showed a consequence for the benchmark's labels.
\par

\noindent\textbf{21.} Counting of the hypotheses. Drafts v2.2 and v2.3 said that twenty-eight entries carried a kill criterion fixed before they ran; H18's pre-registration has none, so there are twenty-seven, of which five survived (draft v1.0 had said that sixteen hypotheses, H6, H13 and H18 among them, were pre-registered each with its kill criterion fixed before execution, and draft v2.1 that all twenty-eight of its entries, T, H6, H13 and H18 included, were written down with a kill criterion). H29 to H31 were counted as pre-registered without the mark that their criterion came from the E11 plan (v2.1 to v2.3), and H2 was not declared as written after its concordance had been observed (all drafts); the table's thirty rows were called thirty-one (v2.2, v2.3); and draft v2.3 attributed the first verification to ``the previous draft'', which was v2.1, not v2.2.
\par

\noindent\textbf{22.} Claims in the abstract and in Section 1. ``Every channel we could construct'' was tested (v2.1 to v2.3), but three activation channels planned in E11 were never run; the comparison with the majority constant did not name its metric (v2.1 to v2.3), and the two representation-distance statistics were scored by AUROC only; the representation distance was listed among ``the statistics that pass their permutation null'' (v2.2, v2.3), while it passes only the uncorrected single-statistic null and not its planned threshold; ``each apparent signal'' was said to trace back to something other than the model's internals (v2.2, v2.3), while our probe keeps about two points over length and the in-sample reading of the artifact was an inference; and the effect of the answer-format prompt was stated without naming Qwen3.5-4B, the only model it concerns (v2.1 to v2.3).
\par

\noindent\textbf{23.} The budget guard. Draft v1.0 (Sections 1 and 8) said it makes timeouts structurally impossible, and drafts v2.1 to v2.3 called it timeout-proof, citing coverage 1.0 in every container evaluation that ran our code. No evaluation came near the limit, the guard never acted, it cannot interrupt a generation under way and relies on a fixed per-item estimate, and several of the evaluations listed did not contain it (Section 1).
\par

\noindent\textbf{24.} H10 (all drafts, Section 5 and PREREG\_E6\_ESITO). The 0.7744 came from a logistic regression on 13 of the 14 official features, three of them constant in our extraction, not from the official regressor and threshold that PREREG\_E6 specified; the substitution was not declared, the outcome file called the protocol executed without changes, and the source of the 0.692 $\pm$ 0.101 it was compared with (the organizers' proposal, a different protocol) was not cited.
\par

\noindent\textbf{25.} H13 (v1.0 \S{}5.2, PREREG\_E8 and the notes added to PREREG\_E7 and E8, and all later drafts in Sections 1, 4 and 5). The oracle's +7.3 was attributed to its freedom to degenerate into a constant. Excluding constant predictions it still gains +5.5; all of its gain comes from thresholds chosen on the items they are scored on, and the fall to $-$2.8 comes from fixing the number of predicted positives, to which Llama contributes nothing.
\par

\noindent\textbf{26.} H14 (all drafts) was described as a realizable calibration via public priors. Each target's public rate is its rate on the tested items, so H14 coincides with the oracle-k ceiling of E8 (both 0.7431); its verdict stands.
\par

\noindent\textbf{27.} Warmup rows changed in draft v2.2 without being declared: H13 moved from dead to a withdrawn ceiling, H16 and H17 from survives to within noise and not established (one of three re-executions of H17's bundle scored 9/17, which meets its kill criterion), H17 had been described as a per-model threshold calibration (only the gpt-oss threshold was calibrated) and H3 as 560 ordered pairs (56 ordered model pairs, 560 predictions), and H6 and T were marked as having no kill criterion.
\par

\noindent\textbf{28.} ``All validation is grouped'' (all drafts). The cross-model tests (H7, H12 to H15) hold out the model but not the problem: 234 of 335 test rows in H7 and 70 of 109 test pairs in H12 to H15 have their problem in training. No verdict changes (Section 4).
\par

\noindent\textbf{29.} H1, H7 and T (all drafts) and H8 (v2.2, v2.3) were computed on the 335 published pairs, counting the 52 duplicated DeepSeek measurements twice. Counted once, H1 is $-$7.4 (not $-$8.7), H8 +1.8 against the global constant (not +1.5), T +7.1 pooled (not +6.0), and H7 stays at +0.0; no verdict changes.
\par

\noindent\textbf{30.} E11 (v2.2, v2.3). The statement that ``the E11 analysis applied Holm to two tests instead of five'' is false: the analysis applied no multiplicity correction, and H30 would have failed Holm over its two tests (permutation p 0.030 against 0.025 at alpha 0.05). Further departures from the plan were not declared: the confirmation set was not held out, the AUROCs were computed on the raw scalars with the orientation chosen on the data, the text baseline of H28 was reduced to character count, the stratified and bookkeeping reference AUROCs that the plan made obligatory were not reported, and only the primary depth was extracted (Section 4).
\par

\noindent\textbf{31.} The V4 tie was presented as caused by re-execution drift in v1.0 (Section 7 and Appendix B) and v2.1 (Appendix B: ``thereby exposed the nondeterminism that Section 8 later explained''), and as ``the first sign of the re-execution noise'' in v2.2 and v2.3 (Section 4), although none of the three explanations recorded at the time was verified; v1.0, v2.2 and v2.3 also quoted its pre-registration without its exception (a tie was impossible unless the DeepSeek path failed). Draft v1.0 (Section 7) and drafts v2.2 and v2.3 (Appendix B) also said that the tie led us to measure re-execution drift; the pre-registration of that measurement (PREREG\_V7) gives as its reason the drift of regressor scores between the logs of v5 and v6 (Appendix B).
\par

\noindent\textbf{32.} The labeling rule. Drafts v1.0 and v2.1 stated that it is exactly recoverable from the labeled rows, and v2.1 that the majority rule of the organizers' proposal agrees with them on only 221 of 379 rows. The labeled rows do not identify the rule (a majority of correct perturbed answers also agrees 379/379), and the 221 is reproduced only by comparing a count of variants with a count of predictions. Draft v2.1 also stated that no public labeled item exists for any of the other four main-phase models, while gpt-oss-120b has ten public pairs, five of them with a declared label, and it called four labeled files aggregated where three are. No draft cited FAQ 7, and none gave the source of the proposal's majority rule (arXiv 2607.13899, its Section 1.4): v2.1 mentioned that rule only through the wrong 221 of 379, and v2.2 and v2.3 dropped the mention instead of correcting it (2.1).
\par

\noindent\textbf{33.} The ground truth (all drafts). The 164 DeepSeek labels were presented as the organizers' labels; 26 of them carry no declared label and were derived by us with the conjunction, and a majority rule, equally consistent with the declared labels, would call 21 of those 26 robust (2.2). Draft v2.3 also stated in 2.2 that the main-phase train set is not used in its analyses, while its 6.4 evaluated our entry on it.
\par

\noindent\textbf{34.} The organizers' sampling regime (v2.1 to v2.3, PREREG\_E13, PREREG\_V8). ``About ten variants per family, one sampled prediction per variant, ten sampled predictions on the base problem'' describes the AIMO-problem file, not the files behind the DeepSeek labels, where most variants were answered ten times and the base between 1 and 60 times; Section 7.2 said the regeneration used ``the organizers' sampling counts'' (v2.2, v2.3) and was titled ``regenerated at their scale'' (v2.1 to v2.3) (2.1, 7.2, 7.3).
\par

\noindent\textbf{35.} Section 2.3(b) of drafts v2.2 and v2.3 gave the official feature's range as 0.959 to 0.999 and the share inside the split range as 91\%, with the thresholds rounded to [0.968, 0.993] as in v2.1; the minimum is 0.958, and between the true lowest and highest thresholds (0.9684 and 0.9935) the share is 92.6\% (98.8\% on DeepSeek).
\par

\noindent\textbf{36.} \texttt{Force\_Latest\_Multiple} (v2.1 to v2.3, OSSERVATO\_MAIN of 2 September, PREREG\_V8, PREREG\_V8b). We read the declared leaderboard rule as keeping only each participant's latest submission per task and justified the resubmission 925822 by it. Under this rule Codabench adds every scored submission as a row of its own and removes none; the valid row 924969 had stayed on the leaderboard, and the resubmission is unaffected (6.3).
\par

\noindent\textbf{37.} V8's parameters. PREREG\_V8 and the bundle's docstring call 98 the median of the character counts of the 174 public problems, and draft v2.1 calls it the median character count; it is the lower median (the median is 98.5). The pre-registration and the docstring give the rule's balanced accuracy on public data as 0.658, the figure for a threshold of 99; drafts v2.1 to v2.3 and a note of 23 September in PREREG\_V8 gave 0.653, itself a double rounding: for the shipped threshold of 98 it is 0.652 over the 335 published pairs and 0.631 over the 283 distinct ones. Drafts v2.2 and v2.3 also said that the rule uses no training data; its threshold and direction were fixed on the public datasets (6.3, Section 13).
\par

\noindent\textbf{38.} Other statements in Section 6. Drafts v2.2 and v2.3 said that whether 11/21 and 10/19 is the majority class ``was not measured'', omitting that on 21 and 19 items the majority constant scores at least 11/21 and 10/19, so neither V8, nor v6, nor the organizers' entries scored above it. Drafts v2.1 to v2.3 described ``a five-model, ten-minute pipeline'', while the runs took 14 minutes with five models and 8.5 minutes with four; and draft v2.3 did not update the scope sentence of 6.1 when it added 6.4.
\par

\noindent\textbf{39.} Section 7.2. The 638 example variants (v2.1 to v2.3) are 624 in the five families (the rest are expert-written or of a sixth family), of which our pools used 621. The statement that on Skywork and Olmo most of the instability comes from which variants are drawn (v2.2, v2.3 and the abstract) holds under the earlier rule for Skywork only: on Olmo the seed alone accounts for 56\% of it. The subsampling curve's concordance ``does not improve'' (v2.1 to v2.3), while it rises from $-$0.02 to +0.03. Draft v2.1 and PREREG\_E13 stated that raising the base-accuracy threshold to 10/10 changes nothing, and drafts v2.2 and v2.3 stated the property behind the unchanged kappas without limiting it to draws A and B. And the table's rate column, draw A over 164 problems, was not labeled as such (v2.1 to v2.3).
\par

\noindent\textbf{40.} The solver and the answer key (v2.1 to v2.3). The ``independent solver'' of 7.2 is gpt-oss-120b itself at a higher reasoning effort, independent of the labeled models but not of the generator; its answer on the original is also the answer key against which every labeled answer was scored, which no draft stated; and the translator of H22 (microsoft/phi-4) was not named.
\par

\noindent\textbf{41.} The simulation of 7.3. Draft v2.1 said that it reproduces the organizers' robust rate, calling 13.8\% of the 116 pairs declared (as does VERIFICA\_classifica\_2026-09-14); it is the share robust under the stricter reading the simulation uses. Draft v2.3 described its data as three per-family files, while one of the three, \texttt{val-sample}, is aggregated and restates the first (721 rows are distinct); it said that its 16 robust pairs were robust ``under the rule'', while 30 are robust under the rule of 2.1 and 16 under that stricter reading; and its abstract and Sections 11 and 12 gave 0.93 where its 7.3 gave 0.94 on average. Drafts v2.1 and v2.2 called the value an upper bound and draft v2.3 called it optimistic, all without the control, run on our own labels, that shows by how much it overstates (about 0.4). No draft said that the simulation covers six model identifiers besides DeepSeek and the organizers' expert families.
\par

\noindent\textbf{42.} Section 7.4 and R3. The headline 80\% (v2.1 to v2.3) came from the 30 problems of the feasibility check; on the full runs the share is 0 and 145 of 164 (88\%). ``DeepSeek-R1 closes spontaneously either way (9.0 of 10)'' (v2.1 to v2.3) gave the bare-text figure for both prompts (9.0 and 9.2). Drafts v2.1 to v2.3 also said that Qwen3.5-4B ``never closes in \textbackslash{}boxed\{\}'' with the bare text, while it did in 5 of the 1,640 base samples of draw A and 8 of the 1,640 of draw B. And R3 said that the prompt, as far as we could find, is not published (v2.3; v2.1 and v2.2 only asked for it to be published), without mentioning the solve prompt in the organizers' baselines repository.
\par

\noindent\textbf{43.} Section 7.5 (v2.2, v2.3). ``Two of our three predictions were wrong'': none held in full, since the third predicted a concordance between 0.10 and 0.40 and it came out at 0.02. The old and new kappas were compared on different sets of problems without saying so, and without their n; on the same problems the change of rule adds 0.01 to 0.11. The gap to the official labels was called one of severity, while our robust rate is the same in both of their classes and our max\_drop does not correlate with theirs. The statement that with the same variants and a new seed agreement is high on Skywork and Olmo, so that the label depends mainly on which variants are drawn, holds for Skywork only: on the 126 problems of Olmo labeled in all three series the seed accounts for about half of the instability. The concordance's upper bound is 0.13, not 0.14; the DeepSeek-identifier rows are 54 on 52 problems; the prediction was 0.60 on three of four models, not on every model; and the abstract's 71 to 75\% did not name the rule it holds under.
\par

\noindent\textbf{44.} Section 8.1 (v2.1 to v2.3). The failed predictions P-E14d and P-E14e were not reported, and the pre-registered measure m4 was not computed; ``each configuration has its own fixed point'' omitted that the two deterministic configurations ran on different cards; crossing counts from different numbers of runs were compared (and R1 set 32 over six runs against 24 over five); and Section 12 named compilation as the inferred cause, which the runs do not single out.
\par

\noindent\textbf{45.} Section 8.2. Its title (``not stable in sign'', v2.1 to v2.3), its statement that one state orders against the rule and the other with it, and the abstract's ``this drift decides whether the official baseline's score channel carries any signal'' (v2.2, v2.3) claimed more than the data show: one state is at chance, the other weakly above it, and their difference is not significant.
\par

\noindent\textbf{46.} Section 8.3. The artifact's stored scores were called out-of-sample scores (v2.1 to v2.3); they are validation balanced accuracies of a rule that almost always predicts one class. Probable in-sample training was offered as the explanation of the artifact's AUROC (v2.1 to v2.3), while its margin tracks problem length and its control probes, trained on shuffled labels, give 0.68 by themselves. ``The probe reads mostly, though not only, problem length'' (v2.2, v2.3) claimed more than the control shows, which is redundancy as predictors. The per-layer maximum is 0.5475, not 0.547 (v2.1 to v2.3); Section 4 said the failed prediction of E15 is reported in 8.3, which did not report it (v2.2, v2.3); and the single-split triple was printed in an order that contradicted its sentence (v2.3).
\par

\noindent\textbf{47.} Section 8.4 (v2.1 to v2.3) did not state that its rates apply the DeepSeek regressor to every model, as our entry does and the official inference does not, nor that gpt-oss-120b is missing because its extraction failed; the introduction of Section 8 called this the unmodified official path.
\par

\noindent\textbf{48.} Compute (v2.1 to v2.3). The Modal total was given as about 27 GPU-hours, a figure no file supports; the dated files document about 12 to 13. The VM total was given as about 100 card-hours, the estimate for E13 alone.
\par

\noindent\textbf{49.} Correspondence. The status of the results was dated 21 September in a draft of 23 September (v2.3), and the note on 8.1 that the E14 pre-registration planned to send by 22 September was not mentioned (v2.2, v2.3); the VM was said to have run Section 8 only (v2.1 to v2.3), while E13 ran there too; E16 was presented as the only experiment on the VM without a kill criterion (v2.2, v2.3), while E13 and E14 had none either; the VM was said to be released ``at the organizers' request'' (v2.2, v2.3), while they asked whether we still needed it; the organizers were said to have raised the cost of averaging three runs in the final evaluation, and our reply to have said that three runs stabilize the published number (v2.1 to v2.3), while their worry concerned running each Codabench evaluation three times during the phase and our reply said ``substantially steadier'' (Section 9); and R4 said the organizers were ``considering'' more rows for Skywork (v2.3), while they wrote that the train set can be expanded.
\par

\noindent\textbf{50.} Record-keeping (v2.2, v2.3). Section 12 of v2.2 and Section 1 of v2.3 said that three analyses had been run from scripts that were not saved; the cross-run analysis of E14, the figures of E16 and of 2.3b, and the balanced accuracy of 6.3 were in the same position. All are now computed by saved scripts (Section 12).
\par

\noindent\textbf{51.} Section 13. The pre-registrations were described as all carrying kill criteria (v2.1 to v2.3; E13, E14 and V7 carry none), and v1.0 (in its Section 11) listed PREREG\_E1 through E8 and v2.1 PREREG\_E1 through E16, although E5, E11 and E16 have no pre-registration file; the deviations declared for the first time omitted E4, E6, most of E11, two of E13 and E14 (v2.2, v2.3), and the list of deviations preserved in place omitted V4 (v2.1 to v2.3); ``no other logged aggregate was used to tune a submission'' omitted H17's motivation and the budget guard's estimate (v2.2, v2.3); the disclosure of AI assistance omitted the bundle code, the VM, the uploads and most of the records it read (v2.2, v2.3); and the release of every dated file did not provide for redacting other participants (v2.1 to v2.3).
\par

\noindent\textbf{52.} Appendix C and Section 6.3. P2 was restated about a single entry (v2.1: ``the entry that led the leaderboard on 7 September''; v2.2, v2.3: ``the entry that led both tracks''), while the dated prediction concerns a participant whose leading entries on the two tracks are distinct submissions. P-V8 was restated without its conditions and P-V8c was omitted (v2.1 to v2.3; v2.1 also omitted the alternative of about 0.66), and Section 6.3 gave only the band of P-V8a, without its condition, as the prediction (v2.1 to v2.3). Two of the three predictions of PREREG\_V8b were omitted (v2.2, v2.3), and the organizers' entries were called baseline entries without a source (v2.2, v2.3).
\par

\noindent\textbf{53.} Section 11's closing sentence (v2.1 to v2.3; v2.1 excepted nothing, v2.2 and v2.3 only P2) did not except P1 and the public leaderboard figures cited in the report, which concern participants' scores as well. Section 9 of drafts v2.2 and v2.3 also contained a statement about other participants' activity in repository issue \#6, which the closing sentence did not except either; it is withdrawn, and Section 9 now says only that the earlier wording was still quoted there on 16 September.
\par

\noindent\textbf{54.} Statements corrected in place without an item until this draft. Draft v1.0 said that the A100 extraction reproduced the organizers' pipeline, with torch and transformers pinned to the container's versions; it was our own re-implementation of the extraction, which differs from the official code on one feature and omits another (2.3b), and torch was 2.13.0, the container's 2.12.1 (Section 5). Draft v2.1 said that every number had been recomputed from disk before being written (several were not, Section 12); that a representation-distance signal at AUROC 0.577 had been exposed as prompt length by the neutral-text control (0.577 is the raw secondary statistic, below its own null; the tested statistic, H30, gave 0.615 and lost to the character count, Section 4); that the gap between the two gpt-oss thresholds was a first sign of the drift of Section 8 (6.2); that the container's main-phase score ranges are higher than the public ones on every model, so that a threshold calibrated on public problems does not transfer (they fall inside the public distributions and cannot tell, 8.4); that Section 8.2 found the 0.7744 of Section 5 not to survive (it was not repeated there, Section 5); that batch size 1 costs 182\% and 73.0 s per item (178\% and 72.3 s on average over its three runs, 8.1); that the second physical card reproduced the first card's state exactly (two of its three runs did, and the third landed in a state of its own, 8.1); that the invalidated Qwen runs took thirteen GPU-hours (about fifteen card-hours, 7.4); and that no other process could regenerate the declared families to useful agreement (7.3, R4). Drafts v2.1 to v2.3 called the regenerated rule the label's stated definition (abstract, Section 7), called our perturbations and sampling variance orthogonal to what the label measures (7.1), described the aggregation as keeping only the detrimental type of a non-robust item, while it keeps the type with the lowest perturbed accuracy and drops milder detrimental ones (2.3a), and described the gpt-oss calibration set as ten public items (v2.1: ``10 public gpt-oss items'') without saying that three of the rows are of a fine-tune and carry its only robust label (6.2). All drafts called \texttt{val-sample} test-matched (abstract, Sections 1 and 4), and drafts v2.2 and v2.3 attributed to the organizers a description of it as covering all test model types, which we have found only in our own dated note (2.4). Drafts v2.2 and v2.3 called the new thresholds the rule the organizers now use and said that the regeneration ran under both of the organizers' labeling rules, while the conjunction is our reading of the data (abstract); said that grouped validation had never been compared with an ungrouped split (Section 4 and item 12), that the organizers' variants break the models more often than ours (only DeepSeek can be compared, 7.3), and that the organizers' labeling prompt and base-accuracy threshold had not been published at the time (they still have not, 7.2; a solve prompt is in their baselines repository, R3); and listed only 8.2 and 8.3 among the measurements that were not pre-registered. Drafts v2.1 to v2.3 also concluded that a learnable target can come only from the organizers (v2.1: ``the only route to a learnable target runs through the organizers' own variants''; v2.2: ``a learnable target has to come from the organizers' own variants''; v2.3: ``the labeled split is the only learnable target''), which our failed regeneration does not show (R4).\par
